% TOP_PROBLEM: {"problemId": "problem.nonvanishing-conjecture-for-projective-log-canonical-pairs", "canonicalTitle": "Nonvanishing conjecture for projective log canonical pairs", "releaseRank": 194, "primaryDomain": "geometry_topology", "primaryDomainLabel": "Geometry & topology", "displayStatus": "open_with_solved_subcases", "releaseStatus": "open_with_solved_subcases"}
% !TeX program = lualatex
% Definition and sources only; this file is not an LLM solution attempt.
\documentclass[11pt]{article}
\usepackage[margin=25mm]{geometry}
\usepackage{fontspec}
\setmainfont{Latin Modern Roman}
\usepackage{amsmath,amssymb,mathtools}
\usepackage{unicode-math}
\setmathfont{Latin Modern Math}
\usepackage{xurl,hyperref}
\hypersetup{colorlinks=true,urlcolor=blue}
\setcounter{secnumdepth}{0}
\setlength{\parindent}{0pt}
\setlength{\parskip}{5pt}
\setlength{\emergencystretch}{3em}
\begin{document}
\section*{\texorpdfstring{Nonvanishing conjecture for projective log canonical pairs}{Nonvanishing conjecture for projective log canonical pairs}}\label{194}
\subsection{Definitions and mathematical statement}
A log canonical pair $(X,\Delta)$ has normal $X$, effective rational boundary $\Delta$, $K_X+\Delta$ rational Cartier, and all discrepancies at least $-1$ on resolutions. A divisor is pseudo-effective if its numerical class lies in the closure of the cone generated by effective divisors. For projective complex pairs, the assertion is
\[
 K_X+\Delta\text{ pseudo-effective}
 \quad\Longrightarrow\quad
 H^0(X,\mathcal O_X(m(K_X+\Delta)))\ne0
\]
for some sufficiently divisible $m>0$ making that divisor integral Cartier. Equivalently its Iitaka dimension is nonnegative. Numerical approximation by effective classes is not itself the existence of a section or linear equivalence to an effective divisor; that is precisely the missing implication.
\par\smallskip\textit{Formulation and concept references:} \href{https://arxiv.org/html/2501.04232v1}{[S1]}.

\subsection{Short English statement}
Must a pseudo-effective log canonical divisor on a projective complex pair acquire a nonzero section after some positive integral multiple?
\subsection{Sources}
\begin{itemize}
\item[S1]Houari Benammar Ammar, Remarks on the Direct Image Sheaf of Logarithmic Pluricanonical Bundles and the Non-Vanishing Conjecture, arXiv:2501.04232v1 (2025).
\url{https://arxiv.org/html/2501.04232v1}.
\textit{Formulation source.}
\end{itemize}
\end{document}
